\bookchapter{Chapitre 34 : Guide approfondi de connectivité automobile — La fusion du smartphone et de la voiture}

\chapterintro{L'automobile passe de « quatre roues et un canapé » à « quatre roues et un smartphone ». Que vous conduisiez une voiture thermique ou à énergie nouvelle, la relation entre smartphone et système embarqué détermine votre expérience de conduite. Ce chapitre clarifie les trois grandes voies de connectivité, les caractéristiques des différents systèmes embarqués, et quelques astuces pour rendre la conduite plus fluide.}

\booksection{I. Les trois voies de connectivité smartphone-voiture}

Il existe actuellement trois approches pour connecter le smartphone et la voiture ; les comprendre, c'est comprendre tout l'écosystème :

\begin{itemize}
  \item \textbf{Duplication d'écran du smartphone vers le système embarqué}: Le système embarqué ne sert que d'« écran » et de « haut-parleur », le calcul et les données cartographiques proviennent du smartphone. Solutions représentatives : Apple CarPlay, Google Android Auto, Huawei HiCar, Baidu CarLife. Avantage : l'écosystème du smartphone est transposé directement dans la voiture, sans réinstaller d'applications ; inconvénient : les capacités matérielles du système embarqué sont sous-utilisées.
  \item \textbf{Intégration profonde smartphone-système embarqué}: Le système embarqué est un système indépendant, mais partage comptes, applications et données avec le smartphone, et peut même « relayer » des tâches. Solutions représentatives : cockpit HarmonyOS (AITO/Luxeed), BYD DiLink, NIO OS, Li Auto. Avantage : expérience intégrée ; inconvénient : généralement lié au même écosystème de marque.
  \item \textbf{CarPlay 2.0 version intégration profonde}: Le nouveau CarPlay lancé par Apple en 2025 ne se contente plus de « dupliquer l'écran », mais prend en charge le tableau de bord, la climatisation, le réglage des sièges et d'autres fonctions centrales du système embarqué. Ferrari, Aston Martin, Porsche et autres marques l'ont adopté en premier, marquant l'entrée des systèmes embarqués dans une nouvelle ère « dirigée par les fabricants de smartphones ».
\end{itemize}

\booksection{II. Comparaison CarPlay / Android Auto / HiCar}

Les trois solutions de duplication dominantes ont chacune leurs atouts ; voici un tableau comparatif pour choisir votre voiture ou votre smartphone :

\begin{center}
\begin{tabularx}{\linewidth}{|X|X|X|X|X|}
\hline
\textbf{Solution} & \textbf{Système smartphone pris en charge} & \textbf{Véhicules pris en charge} & \textbf{Richesse fonctionnelle} & \textbf{Disponibilité en Chine} \tabularnewline
\hline
Apple CarPlay & iPhone uniquement & Très large, la grande majorité des véhicules & Élevée, navigation/musique/téléphone/messages & Bonne, utilisable nativement \tabularnewline
Android Auto & Android (les smartphones nationaux nécessitent le cadre de services Google) & Assez large, surtout véhicules joint-venture/importés & Élevée, comparable à CarPlay & Moyenne, nécessite un VPN pour l'activation \tabularnewline
Huawei HiCar & Surtout Huawei, partiellement autres Android & Véhicules nationaux et certaines joint-ventures & Assez élevée, peut interagir avec le cockpit HarmonyOS & Très bonne, écosystème local \tabularnewline
Baidu CarLife & iOS et Android & Assez large, commun en installation après-vente sur anciens véhicules & Moyenne, mises à jour lentes & Très bonne, écosystème local \tabularnewline
\hline
\end{tabularx}
\end{center}

\booksection{III. Cockpit HarmonyOS : transition fluide entre smartphone et voiture}

Si CarPlay consiste à « projeter le smartphone sur la voiture », le cockpit HarmonyOS considère que « smartphone et voiture ne font qu'un ». La synergie entre les smartphones Huawei et les systèmes embarqués AITO/Luxeed représente le meilleur niveau actuel de connectivité en Chine :

\begin{itemize}
  \item \textbf{Super Bureau}: Le système embarqué appelle directement les applications du smartphone, sans réinstallation. La musique en cours d'écoute ou l'itinéraire de navigation sur le smartphone se poursuivent automatiquement sur l'écran embarqué à l'entrée dans la voiture, et reviennent au smartphone à la sortie, sans interruption.
  \item \textbf{Relais d'application}: La navigation définie sur le smartphone est automatiquement reprise par le système embarqué à la fermeture de la portière ; à l'arrivée au parking, le smartphone reprend la navigation piétonne jusqu'à la destination finale. Le « dernier kilomètre » n'est plus interrompu.
  \item \textbf{Compte partagé smartphone-voiture}: Une fois connecté au compte Huawei, le système embarqué synchronise automatiquement votre abonnement musical, vos favoris cartographiques, votre agenda et vos contacts. La voiture et le smartphone sont différents terminaux d'une même identité, et non deux systèmes indépendants.
\end{itemize}

\booksection{IV. Systèmes embarqués des nouveaux constructeurs}

Les systèmes embarqués des véhicules à énergie nouvelle nationaux sont passés de « navigation + musique » à des terminaux intelligents mobiles, chaque constructeur ayant sa propre approche :

\begin{itemize}
  \item \textbf{NIO OS (NIO)}: L'assistant vocal NOMI est une référence du secteur, capable de reconnaître des commandes multiples et des dialogues contextuels continus, avec des expressions émotionnelles. Le NIO Phone fait du smartphone une « clé de voiture + second écran », profondément lié à l'écosystème NIO.
  \item \textbf{Système embarqué Li Auto}: L'interaction multi-écran est le point fort, l'écran passager peut diffuser des séries indépendamment, l'écran enfant peut lire des dessins animés. Le « Maître des tâches » permet de personnaliser des scénarios, par exemple « pause automatique de la musique à l'ouverture de la portière, reprise à la fermeture ».
  \item \textbf{Système embarqué Xiaomi HyperOS}: Liaison profonde avec les smartphones, montres et appareils électroménagers Xiaomi. L'alarme réglée sur le smartphone rappelle sur le système embarqué ; la navigation vers la maison depuis la voiture déclenche la climatisation à domicile. HyperOS unifie « l'écosystème personne-voiture-maison ».
  \item \textbf{BYD DiLink}: L'écran rotatif emblématique, en mode vertical pour les courtes vidéos, en mode horizontal pour la navigation, s'adapte à l'orientation des applications. Dispose d'une boutique d'applications intégrée, l'écosystème le plus riche parmi les constructeurs traditionnels.
\end{itemize}

\booksection{V. Le smartphone comme système embarqué — est-ce viable ?}

Si votre système embarqué est trop ancien ou trop peu pratique, pouvez-vous utiliser le smartphone à la place ? La réponse est : oui, mais selon les scénarios.

\begin{itemize}
  \item \textbf{Motorola Ready For / Mode bureau}: Certains smartphones Motorola, Samsung et Huawei prennent en charge le « mode bureau », se transformant en ordinateur simplifié connecté à un moniteur. En théorie, cela fonctionne avec un écran embarqué, mais la compatibilité réelle est inégale et il manque les interactions spécifiques à la voiture (boutons au volant).
  \item \textbf{Support de smartphone vs système embarqué d'origine}: Si le système embarqué d'origine est lent ou que les cartes ne sont pas à jour, un support de smartphone solide + navigation sur smartphone est un choix pragmatique. Attention cependant : (1) l'écran du smartphone est petit, difficile à lire en conduisant ; (2) le soleil d'été peut provoquer une surchauffe et un arrêt automatique ; (3) la navigation prolongée consomme beaucoup de batterie, un chargeur de voiture est nécessaire.
  \item \textbf{Recommandation}: Pour les trajets courts ou occasionnels, un support de smartphone suffit ; pour les trajets quotidiens et longue distance, il est recommandé de passer à un système embarqué supportant CarPlay/HiCar ou d'installer un système après-vente, pour une meilleure expérience et sécurité.
\end{itemize}

\begin{tcbwarn}{Rappel de sécurité}
Quelle que soit la solution, manipuler le smartphone en conduisant est une infraction. Le support de smartphone ne sert qu'à « voir la navigation », pas à « manipuler le smartphone ». Pour définir une destination ou changer de musique, arrêtez-vous d'abord ou utilisez le contrôle vocal.

\end{tcbwarn}

\booksection{VI. Astuces pratiques}

\begin{itemize}
  \item \textbf{Connexion CarPlay sans fil instable}: (1) Désactivez puis réactivez le Bluetooth du système embarqué et réappairez ; (2) supprimez l'historique d'appairage du système embarqué et reconnectez ; (3) vérifiez que le smartphone n'est pas connecté à plusieurs appareils Bluetooth simultanément ; (4) mettez à jour le firmware du système embarqué et le système du smartphone ; (5) en dernier recours, utilisez CarPlay filaire, plus stable.
  \item \textbf{Données du système embarqué épuisées}: (1) Activez le point d'accès du smartphone et connectez le système embarqué ; (2) dans les réglages du système embarqué, basculez les applications gourmandes en données (vidéo, musique en ligne) vers le Wi-Fi ; (3) nettoyez régulièrement le cache des cartes hors ligne du système embarqué.
  \item \textbf{Déclenchement automatique du mode conduite}: (1) Le « mode Concentration » d'iOS peut être programmé pour s'activer « à la connexion Bluetooth de la voiture », masquant les notifications non urgentes ; (2) le « mode Conduite » d'Android ou des applications tierces (comme Tasker) peuvent couper le son et répondre automatiquement ; (3) les smartphones Huawei peuvent configurer le « paiement intelligent » pour basculer automatiquement la musique vers le système embarqué.
\end{itemize}

\begin{tcbinfo}{Conseils lors de l'achat d'une voiture}
Si l'expérience du système embarqué compte pour vous, testez en concession : (1) apportez votre smartphone pour tester la stabilité de CarPlay/HiCar ; (2) testez la vitesse de réponse et le taux de reconnaissance de l'assistant vocal ; (3) vérifiez si la boutique d'applications propose vos applications habituelles ; (4) renseignez-vous sur la fréquence des mises à jour OTA — un bon système embarqué s'améliore par itérations logicielles.

\end{tcbinfo}

\begin{tcbpractice}{Exercice : optimisez votre expérience de connectivité voiture}
\begin{enumerate}
  \item Confirmez la solution de connectivité prise en charge par votre smartphone et votre système embarqué (CarPlay / HiCar / Android Auto / CarLife).
  \item Nettoyez une fois la liste d'appairage Bluetooth du système embarqué, ne gardez que les smartphones habituels, pour réduire les conflits de connexion.
  \item Configurez un « mode concentration conduite », s'activant automatiquement à la connexion Bluetooth du système embarqué.
  \item Téléchargez l'application de la marque (comme HarmonyOS Intelligent Driving, Li Auto, NIO), voyez si vous pouvez contrôler la climatisation à distance et consulter l'état du véhicule.
\end{enumerate}

\end{tcbpractice}

\begin{tcbsummary}{Résumé du chapitre}
\begin{itemize}
  \item Trois voies de connectivité : duplication (CarPlay, etc.), intégration profonde (cockpit HarmonyOS, etc.), CarPlay 2.0 intégration profonde
  \item CarPlay a l'écosystème le plus large ; HiCar offre la meilleure expérience en Chine ; Android Auto nécessite des réglages en Chine ; CarLife convient aux anciens véhicules en après-vente
  \item Le Super Bureau, le relais d'application et le compte partagé du cockpit HarmonyOS représentent la direction d'intégration profonde
  \item NOMI de NIO, multi-écran de Li Auto, écosystème personne-voiture-maison de Xiaomi, écran rotatif de BYD : chacun a sa spécificité
  \item Le smartphone comme système embarqué ne convient qu'aux usages temporaires ; pour le long terme, une mise à niveau ou une installation après-vente est recommandée
  \item Astuces : reconnexion CarPlay, point d'accès smartphone pour les données, déclenchement Bluetooth du mode conduite
\end{itemize}

\end{tcbsummary}
